\section{Estrategias de razonamiento: de Few-shot a Zero-shot}

\subsection{Least-to-Most Prompting: descomposición de tareas complejas}

Inspirado en la estrategia de descomposición de la obra clásica de P\'{o}lya \textit{How to Solve It}, el equipo de Denny Zhou propuso Least-to-Most Prompting:

\begin{itemize}
    \item \textbf{Idea central}: primero se descompone el problema complejo en una serie de subproblemas, luego se resuelven de manera progresiva a partir del subproblema más simple, y la solución de cada subproblema sirve como contexto para el siguiente.
    \item \textbf{Tarea SCAN}: punto de referencia de generalización composicional (compositional generalization), donde se alcanza una exactitud de 99.7\% usando solo 1\% de los ejemplos.
    \item \textbf{Tarea Text-to-Code}: usando Dynamic Least-to-Most Prompting, con solo 1\% de los datos se supera ampliamente a arquitecturas especializadas entrenadas con todos los datos de entrenamiento.
\end{itemize}

\begin{knowledgebox}{Generalización composicional (Compositional Generalization)}
Se refiere al escenario en el que las muestras de prueba son más complejas que las muestras de entrenamiento o de los prompts. Por ejemplo, durante el entrenamiento solo se han visto fragmentos de código cortos, pero en la prueba se requiere generar código más largo. Least-to-Most Prompting resuelve este problema de manera natural mediante su estrategia de descomposición.
\end{knowledgebox}

\subsection{Zero-shot CoT: ``Let's think step by step''}

\begin{itemize}
    \item Sin necesidad de ningún ejemplo, basta con añadir la frase ``Let's think step by step'' después de la pregunta para activar el razonamiento paso a paso del LLM.
    \item El desempeño suele ser inferior al de few-shot CoT, pero tiene la ventaja de no tener costo alguno.
\end{itemize}

\subsection{LLMs as Analogical Reasoners}

Inspirado en la idea de razonamiento analógico de P\'{o}lya:

\begin{itemize}
    \item \textbf{Método}: ante un problema nuevo, se hace que el LLM genere por sí mismo problemas relacionados junto con sus soluciones, y luego se utilizan esos ejemplos autogenerados para resolver el problema original.
    \item \textbf{Ventaja}: es notablemente mejor que ``Let's think step by step'', e incluso supera a los ejemplos few-shot diseñados manualmente, porque el modelo genera ejemplos relacionados hechos a la medida de cada problema.
    \item \textbf{Relación con la generación aumentada por recuperación}: puede extenderse además a la recuperación de problemas y conocimientos relacionados desde la red, lo que permite lograr escalabilidad.
\end{itemize}

\subsection{Chain-of-Thought Reasoning without Prompting}

Un hallazgo sorprendente: \textbf{sin necesidad de ningún prompt}, es posible activar el razonamiento únicamente mediante una estrategia de decodificación especial:

\begin{itemize}
    \item En el primer paso de decodificación, en lugar de tomar el token de mayor probabilidad, se exploran los $k$ tokens candidatos principales.
    \item Para cada token candidato se continúa con decodificación voraz, generando una respuesta completa.
    \item Se elige la ruta que contiene pasos de razonamiento y cuya respuesta final tiene la mayor confianza.
\end{itemize}

\begin{importantbox}{Observación clave}
El LLM ya «sabe» internamente cómo razonar. Cuando aparece una ruta de razonamiento (como ``Nicolas Cage was born in 1964, and 1964 is an even year''), la probabilidad de la respuesta final salta de un valor bajo a 98\%. Los pasos de razonamiento aumentan de manera significativa, a nivel de probabilidad, la confianza del modelo en la respuesta final.
\end{importantbox}

\subsection{Resumen de la sección}

De few-shot CoT a zero-shot CoT, y luego a la estrategia de decodificación especial que no requiere prompt alguno, las formas de activar la capacidad de razonamiento del LLM son cada vez más flexibles. La regularidad central es que lo determinante son los pasos de razonamiento en sí mismos, y no la forma de activarlos. El razonamiento analógico muestra además que el LLM puede generar de manera autónoma ejemplos relacionados para apoyar su razonamiento.
